\input{preamble}

% OK, start here.
%
\begin{document}

\title{Quotients de groupoïdes}


\maketitle

\phantomsection
\label{section-phantom}

\tableofcontents

\section{Introduction}
\label{section-introduction}

\noindent
Ce chapitre est consacré à des généralités sur les groupoïdes et leurs
quotients (dans la mesure où ceux-ci existent).
Il existe une abondante littérature sur ce sujet ; voir par exemple
\cite{GIT}, \cite{seshadri_quotients}, \cite{KollarQuotients},
\cite{K-M}, \cite{KollarFinite}, et bien d'autres références.






\section{Conventions et notation}
\label{section-conventions-notation}

\noindent
Dans ce chapitre, les conventions et la notation sont celles introduites dans
Groupoïdes in Espaces, sections \ref{spaces-groupoids-section-conventions}
et \ref{spaces-groupoids-section-notation}.


\section{Morphismes invariants}
\label{section-invariant}

\begin{definition}
\label{definition-invariant}
Soit $S$ un schéma et soit $B$ un espace algébrique sur $S$.
Soit $j = (t, s) : R \to U \times_B U$ une prérelation d'espaces
algébriques sur $B$. On dit qu'un morphisme $\phi : U \to X$ d'espaces algébriques
sur $B$ est {\it $R$-invariant} si le diagramme
$$
\xymatrix{
R \ar[r]_s \ar[d]_t & U \ar[d]^\phi \\
U \ar[r]^\phi & X
}
$$
est commutatif. Si $j : R \to U \times_B U$ provient de l'action
d'un espace algébrique en groupes $G$ sur $U$ au-dessus de $B$, comme dans
Groupoïdes in Espaces, lemme \ref{spaces-groupoids-lemma-groupoid-from-action},
on dit alors que $\phi$ est {\it $G$-invariant}.
\end{definition}

\noindent
Autrement dit, un morphisme $U \to X$ est $R$-invariant s'il égalise
$s$ et $t$. On peut reformuler ceci à l'aide des faisceaux quotients
associés, comme suit.

\begin{lemma}
\label{lemma-invariant}
Soit $S$ un schéma et soit $B$ un espace algébrique sur $S$.
Soit $j = (t, s) : R \to U \times_B U$ une prérelation d'espaces
algébriques sur $B$. Un morphisme d'espaces algébriques $\phi : U \to X$ est
$R$-invariant si et seulement s'il se factorise sous la forme
$U \to U/R \to X$.
\end{lemma}

\begin{proof}
Cela résulte immédiatement de la définition du faisceau quotient donnée dans
Groupoïdes in Espaces, section \ref{spaces-groupoids-section-quotient-sheaves}.
\end{proof}

\begin{lemma}
\label{lemma-base-change-on-invariant}
Soit $S$ un schéma et soit $B$ un espace algébrique sur $S$.
Soit $j = (t, s) : R \to U \times_B U$ une prérelation d'espaces
algébriques sur $B$. Soit $U \to X$ un morphisme $R$-invariant d'espaces
algébriques sur $B$. Soit $X' \to X$ un morphisme quelconque d'espaces algébriques.
\begin{enumerate}
\item En posant $U' = X' \times_X U$, $R' = X' \times_X R$, on obtient
une prérelation $j' : R' \to U' \times_B U'$.
\item Si $j$ est une relation, alors $j'$ est une relation.
\item Si $j$ est une prérelation d'équivalence, alors $j'$ est une
prérelation d'équivalence.
\item Si $j$ est une relation d'équivalence, alors $j'$ est une relation
d'équivalence.
\item Si $j$ provient d'un groupoïde en espaces algébriques
$(U, R, s, t, c)$ sur $B$, alors
\begin{enumerate}
\item $(U, R, s, t, c)$ est un groupoïde en espaces algébriques sur $X$, et
\item $j'$ provient du changement de base $(U', R', s', t', c')$
de ce groupoïde à $X'$ ; voir
Groupoïdes in Espaces, lemme
\ref{spaces-groupoids-lemma-base-change-groupoid}.
\end{enumerate}
\item Si $j$ provient de l'action d'un espace algébrique en groupes $G/B$ sur $U$,
comme dans Groupoïdes in Espaces, lemme
\ref{spaces-groupoids-lemma-groupoid-from-action},
alors $j'$ provient de l'action induite de $G$ sur $U'$.
\end{enumerate}
\end{lemma}

\begin{proof}
Omise. Indication : combiner le point de vue fonctoriel avec le diagramme
$$
\xymatrix{
R' = X' \times_X R \ar[dd] \ar[rr] \ar[rd] & &
X' \times_X U = U' \ar'[d][dd] \ar[rd] \\
& R \ar[dd] \ar[rr] & & U \ar[dd] \\
U' = X' \times_X U \ar'[r][rr] \ar[rd] & & X' \ar[rd] \\
& U \ar[rr] & & X
}
$$
\end{proof}

\begin{definition}
\label{definition-base-change}
Dans la situation du lemme \ref{lemma-base-change-on-invariant},
on appelle $j' : R' \to U' \times_B U'$ le {\it changement de base} de la prérelation
$j$ à $X'$. On dit que c'est un {\it changement de base plat} si $X' \to X$ est un
morphisme plat d'espaces algébriques.
\end{definition}

\noindent
Ce type de changement de base se comporte bien avec la formation des faisceaux
quotients et des champs quotients.

\begin{lemma}
\label{lemma-base-change-quotient-sheaf}
Dans la situation du lemme \ref{lemma-base-change-on-invariant},
il existe un isomorphisme de faisceaux
$$
U'/R' = X' \times_X U/R
$$
Pour la construction des faisceaux quotients, voir
Groupoïdes in Espaces, section \ref{spaces-groupoids-section-quotient-sheaves}.
\end{lemma}

\begin{proof}
Puisque $U \to X$ est $R$-invariant, il est clair que l'application
$U \to X$ se factorise par le faisceau quotient $U/R$.
Rappelons que, par définition,
$$
\xymatrix{
R \ar@<1ex>[r] \ar@<-1ex>[r] &
U \ar[r] &
U/R
}
$$
est un diagramme coégalisateur dans la catégorie $\Sh$ des faisceaux d'ensembles sur
$(\Sch/S)_{fppf}$. En fait, c'est un diagramme coégalisateur dans la
catégorie au-dessus $\Sh/X$. Puisque le foncteur de changement de base
$X' \times_X - : \Sh/X \to \Sh/X'$ est exact (ce qui est vrai dans tout topos),
on conclut.
\end{proof}

\begin{lemma}
\label{lemma-base-change-quotient-stack}
Soit $S$ un schéma. Soit $B$ un espace algébrique sur $S$.
Soit $(U, R, s, t, c)$ un groupoïde en espaces algébriques sur $B$.
Soit $U \to X$ un morphisme $R$-invariant d'espaces algébriques sur
$B$. Soit $g : X' \to X$ un morphisme d'espaces algébriques sur $B$
et soit $(U', R', s', t', c')$ le changement de base défini dans le
lemme \ref{lemma-base-change-on-invariant}. Alors
$$
\xymatrix{
[U'/R'] \ar[r] \ar[d] & [U/R] \ar[d] \\
\mathcal{S}_{X'} \ar[r] & \mathcal{S}_X
}
$$
est un $2$-produit fibré de champs en groupoïdes sur $(\Sch/S)_{fppf}$.
Pour la construction des champs quotients et des morphismes de ce
diagramme, voir
Groupoïdes in Espaces, section \ref{spaces-groupoids-section-stacks}.
\end{lemma}

\begin{proof}
Nous allons le démontrer en utilisant la description explicite
des champs quotients donnée dans
Groupoïdes in Espaces, lemme
\ref{spaces-groupoids-lemma-quotient-stack-objects}.
Nous encourageons néanmoins vivement le lecteur à trouver sa propre démonstration.
Tout d'abord, on peut considérer $(U, R, s, t, c)$ comme un groupoïde en
espaces algébriques sur $X$, d'où une application
$f : [U/R] \to \mathcal{S}_X$ ; voir
Groupoïdes in Espaces, lemme \ref{spaces-groupoids-lemma-quotient-stack-arrows}.
De même, on a $f' : [U'/R'] \to X'$.

\medskip\noindent
Un objet du $2$-produit fibré
$\mathcal{S}_{X'} \times_{\mathcal{S}_X} [U/R]$ au-dessus d'un schéma $T$ sur $S$
revient à un morphisme $x' : T \to X'$ et à un objet $y$ de $[U/R]$ au-dessus de $T$
tels que la composée $g \circ x'$ soit égale à $f(y)$.
Cela a un sens, car les objets de $\mathcal{S}_X$ au-dessus de $T$
sont les morphismes $T \to X$. D'après Groupoïdes in Espaces, lemme
\ref{spaces-groupoids-lemma-quotient-stack-objects},
on peut supposer que $y$ est donné par une donnée de descente $[U/R]$, $(u_i, r_{ij})$,
relative à un recouvrement fppf $\{T_i \to T\}$.
L'égalité $g \circ x' = f(y)$ signifie que les diagrammes
$$
\vcenter{
\xymatrix{
T_i \ar[rr]_{u_i} \ar[d] & & U \ar[d] \\
T \ar[r]^{x'} & X' \ar[r]^g & X
}
}
\quad\text{et}\quad
\vcenter{
\xymatrix{
T_i \times_T T_j \ar[rr]_{r_{ij}} \ar[d] & & R \ar[d] \\
T \ar[r]^{x'} & X' \ar[r]^g & X
}
}
$$
sont commutatifs.

\medskip\noindent
D'autre part, d'après Groupoïdes in Espaces, lemme
\ref{spaces-groupoids-lemma-quotient-stack-objects}, un objet $y'$ de $[U'/R']$
au-dessus d'un schéma $T$ sur $S$
est donné par une donnée de descente $[U'/R']$, $(u'_i, r'_{ij})$,
relative à un recouvrement fppf $\{T_i \to T\}$.
En posant $f'(y') = x' : T \to X'$, on voit que
les diagrammes
$$
\vcenter{
\xymatrix{
T_i \ar[r]_{u'_i} \ar[d] & U' \ar[d] \\
T \ar[r]^{x'} & X'
}
}
\quad\text{et}\quad
\vcenter{
\xymatrix{
T_i \times_T T_j \ar[r]_{r'_{ij}} \ar[d] & U' \ar[d] \\
T \ar[r]^{x'} & X'
}
}
$$
sont commutatifs.

\medskip\noindent
Ces notations étant posées, on définit un foncteur
$$
[U'/R'] \longrightarrow \mathcal{S}_{X'} \times_{\mathcal{S}_X} [U/R]
$$
en envoyant $y' = (u'_i, r'_{ij})$, comme ci-dessus, sur l'objet
$(x', (u_i, r_{ij}))$, où $x' = f'(y')$, où
$u_i$ est la composée $T_i \to U' \to U$, et où
$r_{ij}$ est la composée $T_i \times_T T_j \to R' \to R$.
Réciproquement, étant donné un objet $(x', (u_i, r_{ij})$
du membre de droite, on peut l'envoyer sur l'objet
$((x', u_i), (x', r_{ij}))$ du membre de gauche.
Nous omettons d'indiquer ce qu'il faut faire des morphismes (cela fonctionne
exactement de la même manière).
\end{proof}





\section{Quotients catégoriques}
\label{section-categorical}

\noindent
C'est le type de quotient le plus élémentaire que l'on puisse considérer.

\begin{definition}
\label{definition-categorical}
Soit $S$ un schéma et soit $B$ un espace algébrique sur $S$.
Soit $j = (t, s) : R \to U \times_B U$ une prérelation en espaces algébriques
sur $B$.
\begin{enumerate}
\item On dit qu'un morphisme $\phi : U \to X$ d'espaces algébriques sur $B$
est un {\it quotient catégorique} s'il est $R$-invariant et si,
pour tout morphisme $R$-invariant $\psi : U \to Y$ d'espaces algébriques sur $B$,
il existe un unique morphisme $\chi : X \to Y$ tel que
$\psi = \phi \circ \chi$.
\item Soit $\mathcal{C}$ une sous-catégorie pleine de la catégorie des espaces
algébriques sur $B$. Supposons que $U$, $R$ soient des objets de $\mathcal{C}$.
Dans cette situation, on dit
qu'un morphisme $\phi : U \to X$ d'espaces algébriques sur $B$
est un {\it quotient catégorique dans $\mathcal{C}$}
si $X \in \Ob(\mathcal{C})$, si $\phi$ est $R$-invariant,
et si, pour tout morphisme $R$-invariant
$\psi : U \to Y$ avec $Y \in \Ob(\mathcal{C})$,
il existe un unique morphisme $\chi : X \to Y$ tel
que $\psi = \phi \circ \chi$.
\item Si $B = S$ et si $\mathcal{C}$ est la catégorie des schémas sur $S$,
on dit alors que $U \to X$ est un
{\it quotient catégorique dans la catégorie des schémas}, ou simplement un
{\it quotient catégorique en schémas}.
\end{enumerate}
\end{definition}

\noindent
On distingue souvent une catégorie $\mathcal{C}$ d'espaces algébriques sur $B$
par un axiome de séparation ; voir
l'exemple \ref{example-categories}
pour quelques cas usuels.
Remarquons que $\phi : U \to X$ est un quotient catégorique si et seulement
si $U \to X$ est un coégalisateur des
morphismes $t, s : R \to U$ dans la catégorie. On en déduit immédiatement
le lemme suivant.

\begin{lemma}
\label{lemma-categorical-unique}
Soit $S$ un schéma et soit $B$ un espace algébrique sur $S$.
Soit $j : R \to U \times_B U$ une prérelation en espaces algébriques sur $B$.
S'il existe un quotient catégorique dans la catégorie des espaces algébriques
sur $B$, il est unique à isomorphisme unique près.
Il en va de même pour les quotients catégoriques dans les sous-catégories pleines de
$\textit{Espaces}/B$.
\end{lemma}

\begin{proof}
Voir Catégories, section \ref{categories-section-coequalizers}.
\end{proof}

\begin{example}
\label{example-categories}
Soit $S$ un schéma et soit $B$ un espace algébrique sur $S$.
Voici quelques exemples usuels de catégories $\mathcal{C}$
qui se présentent fréquemment lorsqu'on applique la
définition \ref{definition-categorical} :
\begin{enumerate}
\item $\mathcal{C}$ est la catégorie de tous les espaces algébriques sur $B$ ;
\item $B$ est séparé et $\mathcal{C}$ est la catégorie de tous les espaces
algébriques séparés sur $B$ ;
\item $B$ est quasi-séparé et $\mathcal{C}$ est la catégorie de tous les espaces
algébriques quasi-séparés sur $B$ ;
\item $B$ est localement séparé et $\mathcal{C}$ est la catégorie de tous les espaces
algébriques localement séparés sur $B$ ;
\item $B$ est décent et $\mathcal{C}$ est la catégorie de tous les espaces algébriques
décents sur $B$ ; enfin,
\item $S = B$ et $\mathcal{C}$ est la catégorie des schémas sur $S$.
\end{enumerate}
Dans ces cas, si $\phi : U \to X$ est un quotient catégorique, on dit que
$U \to X$ est
(1) un {\it quotient catégorique},
(2) un {\it quotient catégorique dans les espaces algébriques séparés},
(3) un {\it quotient catégorique dans les espaces algébriques quasi-séparés},
(4) un {\it quotient catégorique dans les espaces algébriques localement séparés},
(5) un {\it quotient catégorique dans les espaces algébriques décents},
(6) un {\it quotient catégorique en schémas}.
\end{example}

\begin{definition}
\label{definition-universal-categorical}
Soit $S$ un schéma et soit $B$ un espace algébrique sur $S$.
Soit $\mathcal{C}$ une sous-catégorie pleine de la catégorie des espaces
algébriques sur $B$, stable par produits fibrés.
Soit $j = (t, s) : R \to U \times_B U$ une prérelation dans
$\mathcal{C}$, et soit $U \to X$ un morphisme $R$-invariant avec
$X \in \Ob(\mathcal{C})$.
\begin{enumerate}
\item On dit que $U \to X$ est un {\it quotient catégorique universel}
dans $\mathcal{C}$ si, pour tout morphisme $X' \to X$ dans $\mathcal{C}$,
le morphisme $U' = X' \times_X U \to X'$ est le quotient catégorique dans
$\mathcal{C}$ du changement de base $j' : R' \to U'$ de $j$.
\item On dit que $U \to X$ est un {\it quotient catégorique uniforme}
dans $\mathcal{C}$ si, pour tout morphisme plat $X' \to X$ dans $\mathcal{C}$,
le morphisme $U' = X' \times_X U \to X'$ est le quotient catégorique dans
$\mathcal{C}$ du changement de base $j' : R' \to U'$ de $j$.
\end{enumerate}
\end{definition}

\begin{lemma}
\label{lemma-categorical-reduced}
Dans la situation de la
définition \ref{definition-categorical},
si $\phi : U \to X$ est un quotient catégorique et si $U$ est réduit,
alors $X$ est réduit. Le même énoncé vaut pour les quotients catégoriques dans
les catégories d'espaces $\mathcal{C}$ énumérées dans
l'exemple \ref{example-categories}.
\end{lemma}

\begin{proof}
Soit $X_{red}$ la réduction de l'espace algébrique $X$.
Puisque $U$ est réduit, le morphisme $\phi : U \to X$ se factorise par
$i : X_{red} \to X$ (Propriétés des espaces algébriques, lemme
\ref{spaces-properties-lemma-map-into-reduction}). Notons ce morphisme
$\phi_{red} : U \to X_{red}$. Comme $\phi \circ s = \phi \circ t$, on
voit aussi que $\phi_{red} \circ s = \phi_{red} \circ t$ (car
$i : X_{red} \to X$ est un monomorphisme). La propriété universelle
de $\phi$ donne donc un morphisme $\chi : X \to X_{red}$ tel que
$\phi_{red} = \phi \circ \chi$. Par unicité, on voit que
$i \circ \chi = \text{id}_X$ et $\chi \circ i = \text{id}_{X_{red}}$.
Ainsi, $i$ est un isomorphisme et $X$ est réduit.

\medskip\noindent
Pour voir que cet argument fonctionne dans une catégorie $\mathcal{C}$, il
suffit de montrer que la réduction d'un objet de $\mathcal{C}$
est un objet de $\mathcal{C}$. Nous omettons de vérifier que c'est
le cas pour chacun des exemples usuels.
\end{proof}





\section{Quotients comme espaces d'orbites}
\label{section-orbits}

\noindent
Soit $j = (t, s) : R \to U \times_B U$ une prérelation.
Si $j$ est une prérelation d'équivalence, alors, de façon informelle,
les « orbites » de $R$ dans $U$
sont les sous-ensembles $t(s^{-1}(\{u\}))$ de $U$. Cependant, si $j$ n'est qu'une
prérelation, il faut prendre la relation d'équivalence engendrée
par $R$.

\begin{definition}
\label{definition-orbit}
Soit $S$ un schéma et soit $B$ un espace algébrique sur $S$.
Soit $j : R \to U \times_B U$ une prérelation sur $B$.
Si $u \in |U|$, alors l'{\it orbite}, ou plus précisément l'
{\it $R$-orbite} de $u$, est
$$
O_u =
\left\{
u' \in |U|\ :
\begin{matrix}
\exists n \geq 1, \ \exists u_0, \ldots, u_n \in |U|\text{ tels que }
u_0 = u \text{ et } u_n = u' \\
\text{et, pour tout }i \in \{0, \ldots, n - 1\},\text{ soit }
u_i = u_{i + 1}\text{ or } \\
\exists r \in |R|, \ s(r) = u_i, t(r) = u_{i + 1}
\text{ or } \\
\exists r \in |R|, \ t(r) = u_i, s(r) = u_{i + 1}
\end{matrix}
\right\}
$$
\end{definition}

\noindent
Il est clair que ce sont les classes d'équivalence d'une relation d'équivalence,
c'est-à-dire que $u' \in O_u$ si et seulement si $u \in O_{u'}$. Le lemme suivant
est une reformulation de
Groupoïdes in Espaces,
lemme \ref{spaces-groupoids-lemma-pre-equivalence-equivalence-relation-points}.

\begin{lemma}
\label{lemma-pre-equivalence-equivalence-relation-points}
Soit $B \to S$ comme dans la section \ref{section-conventions-notation}.
Soit $j : R \to U \times_B U$ une prérelation d'équivalence
d'espaces algébriques sur $B$. Alors
$$
O_u =
\{u' \in |U| \text{ tel qu'il existe } r \in |R|, \ s(r) = u, \ t(r) = u'\}.
$$
\end{lemma}

\begin{proof}
D'après le lemme précité de
Groupoïdes in Espaces,
lemme \ref{spaces-groupoids-lemma-pre-equivalence-equivalence-relation-points},
les orbites $O_u$ définies dans le lemme donnent une décomposition de
$|U|$ en réunion disjointe. Elles sont donc égales aux
orbites définies dans la définition \ref{definition-orbit}.
\end{proof}

\begin{lemma}
\label{lemma-invariant-map-constant-on-orbit}
Dans la situation de la définition \ref{definition-orbit},
soit $\phi : U \to X$ un morphisme $R$-invariant d'espaces algébriques sur
$B$. Alors $|\phi| : |U| \to |X|$ est constant sur les orbites.
\end{lemma}

\begin{proof}
Il suffit de montrer que $\phi(u) = \phi(u')$
pour tous $u, u' \in |U|$ tels qu'il existe
$r \in |R|$ vérifiant $s(r) = u$ et $t(r) = u'$.
Cela est clair puisque $\phi$ égalise $s$ et $t$.
\end{proof}

\noindent
L'utilisation des orbites $O_u \subset |U|$ pour caractériser les propriétés
des applications quotients présente plusieurs difficultés. En voici une.
Supposons que $\Spec(k) \to B$
soit un point géométrique de $B$. Considérons l'application canonique
$$
U(k) \longrightarrow |U|.
$$
En général, les classes d'équivalence
de la relation d'équivalence engendrée par $j(R(k)) \subset U(k) \times U(k)$
ne sont pas les images réciproques des orbites $O_u \subset |U|$.
Un exemple élémentaire consiste à prendre $S = B = \Spec(\mathbf{Z})$,
$U = R = \Spec(k)$ avec $s = t = \text{id}_k$. Alors $|U| = |R|$ se réduit
à un point, mais $U(k)/R(k)$ est énorme.
Un exemple plus intéressant consiste à prendre $S = B = \Spec(\mathbf{Q})$,
à choisir des corps de nombres $K \subset L$, puis à poser $U = \Spec(L)$
et $R = \Spec(L \otimes_K L)$, avec les applications évidentes $s, t : R \to U$.
Dans ce cas, $|U|$ n'a toujours qu'un point, mais le quotient
$$
U(k)/R(k) = \Hom(K, k)
$$
contient plus d'un élément. Ces deux exemples montrent
que, si $U \to X$ est une application $R$-invariante et si l'on veut qu'elle
« sépare les orbites », on obtient une notion beaucoup plus forte et intéressante en
considérant les applications induites $U(k) \to X(k)$ et en exigeant que
ces applications séparent les orbites.

\medskip\noindent
Cette formulation présente encore une difficulté. Supposons en effet que
$S = B = \Spec(\mathbf{R})$,
$U = \Spec(\mathbf{C})$, et
$R = \Spec(\mathbf{C}) \amalg \Spec(K)$
pour une extension de corps $\sigma : \mathbf{C} \to K$.
Les applications $s, t$ sont données par l'identité sur la composante $\Spec(\mathbf{C})$,
mais, sur la seconde composante, elles sont données par $\sigma, \sigma \circ \tau$, où
$\tau$ est la conjugaison complexe. Si
$K$ est une extension non triviale de $\mathbf{C}$, les deux points
$1, \tau \in U(\mathbf{C})$ ne sont pas équivalents modulo
$j(R(\mathbf{C}))$. Toutefois, après avoir choisi une extension $\mathbf{C} \subset \Omega$
de cardinal suffisamment grand (par exemple supérieur au cardinal
de $K$), les images de $1, \tau \in U(\mathbf{C})$ dans
$U(\Omega)$ deviennent équivalentes ! Intuitivement, il semble clair que
cela se produit soit parce que $s, t : R \to U$ ne sont pas localement de type fini,
soit parce que le cardinal du corps $k$ n'est pas assez grand.

\medskip\noindent
Compte tenu de ce qui précède, posons la définition suivante.

\begin{definition}
\label{definition-geometric-orbits}
Soit $S$ un schéma et soit $B$ un espace algébrique sur $S$.
Soit $j : R \to U \times_B U$ une prérelation sur $B$.
Soit $\Spec(k) \to B$ un point géométrique de $B$.
\begin{enumerate}
\item On dit que $\overline{u}, \overline{u}' \in U(k)$ sont
{\it faiblement $R$-équivalents} s'ils appartiennent à la même classe d'équivalence
pour la relation d'équivalence engendrée par la relation
$j(R(k)) \subset U(k) \times U(k)$.
\item On dit que $\overline{u}, \overline{u}' \in U(k)$ sont
{\it $R$-équivalents} si, pour un surcorps $k \subset \Omega$,
leurs images dans $U(\Omega)$ sont faiblement $R$-équivalentes.
\item L'{\it orbite faible}, ou plus précisément l'{\it $R$-orbite faible},
de $\overline{u} \in U(k)$ est l'ensemble de tous les
éléments de $U(k)$ faiblement $R$-équivalents à $\overline{u}$.
\item L'{\it orbite}, ou plus précisément l'{\it $R$-orbite},
de $\overline{u} \in U(k)$ est l'ensemble de tous les
éléments de $U(k)$ qui sont $R$-équivalents à $\overline{u}$.
\end{enumerate}
\end{definition}

\noindent
On verra que, dans les cas favorables, orbites et orbites faibles coïncident ; voir
le lemme \ref{lemma-geometric-orbits}. Le lemme suivant illustre
la différence dans le cas particulier d'une prérelation d'équivalence.

\begin{lemma}
\label{lemma-weak-orbit-pre-equivalence}
Soit $S$ un schéma et soit $B$ un espace algébrique sur $S$.
Soit $\Spec(k) \to B$ un point géométrique de $B$.
Soit $j : R \to U \times_B U$ une prérelation d'équivalence sur $B$.
Dans ce cas, l'orbite faible de $\overline{u} \in U(k)$ est simplement
$$
\{
\overline{u}' \in U(k)
\text{ tel que }
\exists \overline{r} \in R(k),
\ s(\overline{r}) = \overline{u},
\ t(\overline{r}) = \overline{u}'
\}
$$
et l'orbite de $\overline{u} \in U(k)$ est
$$
\{
\overline{u}' \in U(k) :
\exists\text{ field extension }K/k, \ \exists\ r \in R(K),
\ s(r) = \overline{u}, \ t(r) = \overline{u}'\}
$$
\end{lemma}

\begin{proof}
Cela résulte du fait que, par définition d'une prérelation d'équivalence, l'image
$j(R(k)) \subset U(k) \times U(k)$ est une relation d'équivalence.
\end{proof}

\noindent
Décrivons le procédé qui transforme toute prérelation en une
prérelation d'équivalence. Nous utiliserons les morphismes
\begin{equation}
\label{equation-list}
\begin{matrix}
j_{diag} &
: &
U &
\longrightarrow &
U \times_B U, &
u &
\longmapsto &
(u, u) \\
j_{flip} &
: &
R &
\longrightarrow &
U \times_B U, &
r &
\longmapsto &
(s(r), t(r)) \\
j_{comp} &
: &
R \times_{s, U, t} R &
\longrightarrow &
U \times_B U, &
(r, r') &
\longmapsto &
(t(r), s(r'))
\end{matrix}
\end{equation}
On définit $j_1 = (t_1, s_1) : R_1 \to U \times_B U$ comme étant le morphisme
$$
j \amalg j_{diag} \amalg j_{flip} :
R \amalg U \amalg R
\longrightarrow
U \times_B U
$$
avec les notations de
l'équation (\ref{equation-list}).
Pour $n > 1$, on pose
$$
j_n = (t_n, s_n) :
R_n = R_1 \times_{s_1, U, t_{n - 1}} R_{n - 1} \longrightarrow U \times_B U
$$
où $t_n$ provient de $t_1$ précomposé avec la projection sur $R_1$, et
$s_n$ de $s_{n - 1}$ précomposé avec la projection sur $R_{n - 1}$.
Enfin, on note
$$
j_\infty = (t_\infty, s_\infty) :
R_\infty = \coprod\nolimits_{n \geq 1} R_n
\longrightarrow
U \times_B U.
$$

\begin{lemma}
\label{lemma-make-pre-equivalence}
Soit $S$ un schéma et soit $B$ un espace algébrique sur $S$.
Soit $j : R \to U \times_B U$ une prérelation sur $B$.
Alors $j_\infty : R_\infty \to U \times_B U$ est une
prérelation d'équivalence sur $B$. De plus,
\begin{enumerate}
\item $\phi : U \to X$ est $R$-invariant si et seulement s'il est
$R_\infty$-invariant ;
\item l'application canonique de faisceaux quotients $U/R \to U/R_\infty$ (voir
Groupoïdes in Espaces, section \ref{spaces-groupoids-section-quotient-sheaves})
est un isomorphisme ;
\item les $R$-orbites faibles coïncident avec les $R_\infty$-orbites faibles ;
\item les $R$-orbites coïncident avec les $R_\infty$-orbites ;
\item si $s, t$ sont localement de type fini, alors $s_\infty$, $t_\infty$
sont localement de type fini ;
\item ajouter ici d'autres propriétés au besoin.
\end{enumerate}
\end{lemma}

\begin{proof}
Omise. Indication pour (5) : toute propriété de $s, t$ stable par composition,
stable par changement de base et locale pour la topologie de Zariski sur la source
est héritée par $s_\infty, t_\infty$.
\end{proof}

\begin{lemma}
\label{lemma-geometric-orbits}
Soit $S$ un schéma et soit $B$ un espace algébrique sur $S$.
Soit $j : R \to U \times_B U$ une prérelation sur $B$.
Soit $\Spec(k) \to B$ un point géométrique de $B$.
\begin{enumerate}
\item Si $s, t : R \to U$ sont localement de type fini,
alors l'équivalence $R$-faible sur $U(k)$ coïncide avec l'$R$-équivalence, et
les $R$-orbites faibles coïncident avec les $R$-orbites dans $U(k)$.
\item Si $k$ est de cardinal suffisamment grand, alors l'équivalence $R$-faible
sur $U(k)$ coïncide avec l'$R$-équivalence, et les $R$-orbites faibles
coïncident avec les $R$-orbites dans $U(k)$.
\end{enumerate}
\end{lemma}

\begin{proof}
Démontrons d'abord (1). Supposons $s, t$ localement de type fini. D'après le
lemme \ref{lemma-make-pre-equivalence},
on peut supposer que $R$ est une prérelation d'équivalence.
Soit $k$ un corps algébriquement clos au-dessus de $B$.
Supposons que $\overline{u}, \overline{u}' \in U(k)$ soient $R$-équivalents.
Alors, pour une extension de corps $\Omega/k$, il existe
un point $\overline{r} \in R(\Omega)$ qui s'envoie sur
$(\overline{u}, \overline{u}') \in (U \times_B U)(\Omega)$ ; voir le
lemme \ref{lemma-weak-orbit-pre-equivalence}.
Ainsi,
$$
Z = R \times_{j, U \times_B U, (\overline{u}, \overline{u}')} \Spec(k)
$$
est non vide. Comme $s$ est localement de type fini,
$j$ l'est aussi ; voir
Morphismes d'espaces algébriques, lemme \ref{spaces-morphisms-lemma-permanence-finite-type}.
Il s'ensuit que $Z$ est un espace algébrique non vide, localement de type fini
sur le corps algébriquement clos $k$ (utiliser
Morphismes d'espaces algébriques,
lemme \ref{spaces-morphisms-lemma-base-change-finite-type}).
Par conséquent, $Z$ possède un point à valeurs dans $k$ ; voir
Morphismes d'espaces algébriques, lemme
\ref{spaces-morphisms-lemma-locally-finite-type-surjective-geometric-points}.
Il existe donc $\overline{r} \in R(k)$ tel que
$j(\overline{r}) = (\overline{u}, \overline{u}')$, et l'on conclut comme voulu que
$\overline{u}, \overline{u}'$ sont $R$-équivalents.

\medskip\noindent
La démonstration de (2) est la même, si ce n'est qu'elle utilise
Morphismes d'espaces algébriques, lemme
\ref{spaces-morphisms-lemma-large-enough}
au lieu de
Morphismes d'espaces algébriques, lemme
\ref{spaces-morphisms-lemma-locally-finite-type-surjective-geometric-points}.
Cela montre que l'assertion vaut dès que $|k| > \lambda(R)$, où
$\lambda(R)$ est introduit juste avant
Morphismes d'espaces algébriques, lemme
\ref{spaces-morphisms-lemma-locally-finite-type-surjective-geometric-points}.
\end{proof}

\noindent
Dans la définition suivante, dire que « $k$ est un corps
au-dessus de $B$ » signifie que $\Spec(k)$ est muni d'un morphisme
$\Spec(k) \to B$.

\begin{definition}
\label{definition-set-theoretically-invariant}
Soit $S$ un schéma et soit $B$ un espace algébrique sur $S$.
Soit $j : R \to U \times_B U$ une prérelation sur $B$.
\begin{enumerate}
\item On dit que $\phi : U \to X$ est {\it ensemblistement $R$-invariant}
si et seulement si l'application $U(k) \to X(k)$ égalise les deux applications
$s, t : R(k) \to U(k)$ pour tout corps algébriquement clos $k$
au-dessus de $B$.
\item On dit que $\phi : U \to X$ {\it sépare les orbites}, ou
{\it sépare les $R$-orbites}, s'il est ensemblistement $R$-invariant et si
$\phi(\overline{u}) = \phi(\overline{u}')$ dans $X(k)$ implique que
$\overline{u}, \overline{u}' \in U(k)$ appartiennent à la même orbite,
pour tout corps algébriquement clos $k$ au-dessus de $B$.
\end{enumerate}
\end{definition}

\noindent
Dans
l'exemple \ref{example-not-invariant},
on montre que l'invariance ensembliste est une notion « trop faible » dans
la catégorie des espaces algébriques. Une reformulation plus géométrique
de l'invariance ensembliste et de la séparation des orbites est donnée au
lemme \ref{lemma-separates-orbits}.

\begin{lemma}
\label{lemma-set-theoretic-invariant}
Dans la situation de la définition \ref{definition-set-theoretically-invariant},
un morphisme $\phi : U \to X$ est ensemblistement $R$-invariant si et
seulement si, pour tout corps algébriquement clos $k$ au-dessus de $B$, l'application
$U(k) \to X(k)$ est constante sur les orbites.
\end{lemma}

\begin{proof}
Cela résulte du fait que la condition est imposée pour tous les corps
algébriquement clos au-dessus de $B$.
\end{proof}

\begin{lemma}
\label{lemma-invariant-set-theoretically-invariant}
Dans la situation de la définition \ref{definition-set-theoretically-invariant},
un morphisme invariant est ensemblistement invariant.
\end{lemma}

\begin{proof}
Cela découle immédiatement des définitions.
\end{proof}

\begin{lemma}
\label{lemma-set-theoretically-invariant-invariant-when-reduced}
Dans la situation de la définition \ref{definition-set-theoretically-invariant},
soit $\phi : U \to X$ un morphisme d'espaces algébriques sur $B$.
Supposons que
\begin{enumerate}
\item $\phi$ soit ensemblistement $R$-invariant ;
\item $R$ soit réduit ; et
\item $X$ soit localement séparé sur $B$.
\end{enumerate}
Alors $\phi$ est $R$-invariant.
\end{lemma}

\begin{proof}
Considérons l'égalisateur
$$
Z = R \times_{(\phi, \phi) \circ j, X \times_B X, \Delta_{X/B}} X
$$
comme espace algébrique. Alors $Z \to R$ est une immersion en vertu de (3).
D'après (1), $|Z| \to |R|$ est surjectif. Il s'ensuit que
$Z \to R$ est une immersion fermée bijective (utiliser
Schémas, lemme \ref{schemes-lemma-immersion-when-closed})
et, par (2), on conclut que $Z = R$.
\end{proof}

\begin{example}
\label{example-not-invariant}
Il existe des espaces algébriques réduits quasi-séparés $X$, $Y$ et deux
morphismes $a, b : Y \to X$ qui coïncident sur tous les points à valeurs dans $k$ sans être
égaux. Pour obtenir un exemple, prendre $Y = \Spec(k[[x]])$ et
$$
X = \mathbf{A}^1_k \Big/ \big(\Delta \amalg \{(x, -x) \mid x \not = 0\}\big)
$$
l'espace algébrique de
Espaces, exemple \ref{spaces-example-affine-line-involution}.
Les deux morphismes $a, b : Y \to X$
proviennent des deux applications $x \mapsto x$ et $x \mapsto -x$
de $Y$ vers $\mathbf{A}^1_k = \Spec(k[x])$. Au point générique,
les deux applications sont égales car, sur l'ouvert $x \not = 0$ de l'espace
$X$, les fonctions $x$ et $-x$ sont égales. Au point fermé,
les applications sont manifestement égales. On a aussi $a \not = b$.
Cela implique que le
lemme \ref{lemma-set-theoretically-invariant-invariant-when-reduced}
cesse d'être vrai si l'on remplace (3) par l'hypothèse que $X$
est quasi-séparé. En effet, considérons le diagramme
$$
\xymatrix{
Y \ar[d]_{-1} \ar[r]_1 & Y \ar[d]^a \\
Y \ar[r]^a & X
}
$$
alors la composée $a \circ (-1) = b$. On peut donc poser $R = Y$,
$U = Y$, $s = 1$, $t = -1$, $\phi = a$ pour obtenir un exemple de morphisme
ensemblistement invariant qui n'est pas invariant.
\end{example}

\noindent
L'exemple précédent est instructif car l'application $Y \to X$ sépare même
les orbites. Il montre que, dans la catégorie des espaces algébriques, il existe tout
simplement trop de morphismes ensemblistement invariants. Définissons maintenant
ce que signifie pour $R$ d'être une relation d'équivalence ensembliste, en gardant
à l'esprit qu'il faut autoriser des extensions de corps pour que la définition
soit correcte.

\begin{definition}
\label{definition-set-theoretic-equivalence}
Soit $S$ un schéma et soit $B$ un espace algébrique sur $S$.
Soit $j : R \to U \times_B U$ une prérelation sur $B$.
\begin{enumerate}
\item On dit que $j$ est une {\it prérelation d'équivalence ensembliste} si,
pour tout corps algébriquement clos $k$ au-dessus de $B$, la relation
$\sim_R$ sur $U(k)$ définie par
$$
\overline{u} \sim_R \overline{u}'
\Leftrightarrow
\begin{matrix}
\exists\text{ field extension }K/k, \ \exists\ r \in R(K), \\
s(r) = \overline{u}, \ t(r) = \overline{u}'
\end{matrix}
$$
est une relation d'équivalence.
\item On dit que $j$ est une {\it relation d'équivalence ensembliste}
si $j$ est universellement injectif et est une prérelation d'équivalence
ensembliste.
\end{enumerate}
\end{definition}

\noindent
Reformulons cette définition en termes plus géométriques.

\begin{lemma}
\label{lemma-set-theoretic-pre-equivalence-geometric}
Dans la situation de la définition \ref{definition-set-theoretic-equivalence},
les conditions suivantes sont équivalentes :
\begin{enumerate}
\item Le morphisme $j$ est une prérelation d'équivalence ensembliste.
\item Le sous-ensemble $j(|R|) \subset |U \times_B U|$ contient l'image de
$|j'|$ pour chacun des morphismes $j'$ de l'équation (\ref{equation-list}).
\item Pour tout corps algébriquement clos $k$ au-dessus de $B$ et de cardinal suffisamment grand,
le sous-ensemble $j(R(k)) \subset U(k) \times U(k)$ est une relation
d'équivalence.
\end{enumerate}
Si $s, t$ sont localement de type fini, ces conditions équivalent aussi à
\begin{enumerate}
\item[(4)] Pour tout corps algébriquement clos $k$ au-dessus de $B$,
le sous-ensemble $j(R(k)) \subset U(k) \times U(k)$ est une relation d'équivalence.
\end{enumerate}
\end{lemma}

\begin{proof}
Supposons (2). Soit $k$ un corps algébriquement clos au-dessus de $B$.
Montrons que $\sim_R$ est une relation d'équivalence.
Supposons que $\overline{u}_i : \Spec(k) \to U$, $i = 1, 2$,
soient des points à valeurs dans $k$ de $U$. Supposons que $(\overline{u}_1, \overline{u}_2)$
soit l'image d'un point à valeurs dans $K$, $r \in R(K)$. Considérons le
diagramme commutatif en traits pleins
$$
\xymatrix{
\Spec(K') \ar@{..>}[r] \ar@{..>}[d]
&
\Spec(k) \ar[d]_{(\overline{u}_2, \overline{u}_1)} &
\Spec(K) \ar[d] \ar[l] \\
R \ar[r]^-j &
U \times_B U &
R \ar[l]_-{j_{flip}}
}
$$
On note aussi $r \in |R|$ l'image de $r$.
Par hypothèse, l'image de $|j_{flip}|$ est contenue dans l'image de
$|j|$ ; autrement dit, il existe $r' \in |R|$ tel que
$|j|(r') = |j_{flip}|(r)$. Mais remarquons que $(\overline{u}_2, \overline{u}_1)$
appartient à la classe d'équivalence qui définit $|j|(r')$ (par commutativité
de la partie en traits pleins du diagramme). Il existe donc une extension de corps
$K'/k$ et un morphisme $r' : \Spec(K) \to R$
(encore noté abusivement $r'$) tels que
$j \circ r' = (\overline{u}_2, \overline{u}_1) \circ i$,
où $i : \Spec(K') \to \Spec(K)$ est l'application évidente.
Autrement dit, la partie pointillée du diagramme est commutative.
Cela prouve que $\sim_R$ est une relation symétrique sur $U(k)$.
De même, le fait que l'image de $|j_{diag}|$ soit contenue
dans l'image de $|j|$ montre que $\sim_R$ est réflexive (détails omis).

\medskip\noindent
Pour montrer que $\sim_R$ est transitive, supposons donnés
$\overline{u}_i : \Spec(k) \to U$, $i = 1, 2, 3$,
des extensions de corps $K_i/k$ et des points
$r_i : \Spec(K_i) \to R$, $i = 1, 2$, tels que
$j(r_1) = (\overline{u}_1, \overline{u}_2)$ et
$j(r_1) = (\overline{u}_2, \overline{u}_3)$. On peut alors choisir un
diagramme commutatif de corps
$$
\xymatrix{
K & K_2 \ar[l] \\
K_1 \ar[u] & k \ar[l] \ar[u]
}
$$
et considérer $r_1, r_2 \in R(K)$. Considérons le
diagramme commutatif en traits pleins
$$
\xymatrix{
\Spec(K') \ar@{..>}[r] \ar@{..>}[d]
&
\Spec(k) \ar[d]_{(\overline{u}_1, \overline{u}_3)} &
\Spec(K) \ar[d]^{(r_1, r_2)} \ar[l]
\\
R \ar[r]^-j &
U \times_B U &
R \times_{s, U, t} R \ar[l]_-{j_{comp}}
}
$$
Le même raisonnement que dans la première partie de la démonstration, en utilisant
cette fois que $|j_{comp}|((r_1, r_2))$ appartient à l'image de $|j|$,
montre qu'il existe un corps $K'$ et des flèches pointillées qui rendent
le diagramme commutatif. Cela prouve que $\sim_R$ est transitive et achève
la preuve de l'implication (2) vers (1).

\medskip\noindent
Supposons (1) et soit $k$ un corps algébriquement clos au-dessus de $B$, dont le
cardinal est supérieur à $\lambda(R)$ ; voir
Morphismes d'espaces algébriques, lemme \ref{spaces-morphisms-lemma-large-enough}.
Supposons que $\overline{u} \sim_R \overline{u}'$, avec
$\overline{u}, \overline{u}' \in U(k)$. Par hypothèse, il existe
un point de $|R|$ qui s'envoie sur $(\overline{u}, \overline{u}') \in |U \times_B U|$.
Ainsi, d'après
Morphismes d'espaces algébriques, lemme \ref{spaces-morphisms-lemma-large-enough},
il existe $\overline{r} \in R(k)$ tel que
$j(\overline{r}) = (\overline{u}, \overline{u}')$. On voit ainsi
que (1) implique (3).

\medskip\noindent
Supposons (3). Montrons que
$\Im(|j_{comp}|) \subset \Im(|j|)$. Choisissons un point quelconque
$c \in |R \times_{s, U, t} R|$. On peut le représenter par un morphisme
$\overline{c} : \Spec(k) \to R \times_{s, U, t} R$, le corps $k$ au-dessus de $B$
étant de cardinal suffisamment grand. Par hypothèse,
$j_{comp}(\overline{c}) \in U(k) \times U(k) = (U \times_B U)(k)$
est aussi l'image $j(\overline{r})$ d'un certain $\overline{r} \in R(k)$.
Par conséquent, $j_{comp}(c) = j(r)$ dans $|U \times_B U|$, comme voulu (où
$r \in |R|$ est la classe d'équivalence de $\overline{r}$). Le même argument
montre aussi que $\Im(|j_{diag}|) \subset \Im(|j|)$ et
$\Im(|j_{flip}|) \subset \Im(|j|)$ (détails omis).
On voit ainsi que (3) implique (2). Nous avons donc montré que
(1), (2) et (3) sont toutes équivalentes.

\medskip\noindent
Il est clair que (4) implique (3) (sans aucune hypothèse sur $s$, $t$).
Pour achever la démonstration du lemme, montrons que (1) implique (4) si $s, t$
sont localement de type fini. Soit donc $k$ un corps algébriquement clos
au-dessus de $B$. Supposons que $\overline{u} \sim_R \overline{u}'$, avec
$\overline{u}, \overline{u}' \in U(k)$. Par hypothèse, l'espace algébrique
$Z = R \times_{j, U \times_B U, (\overline{u}, \overline{u}')} \Spec(k)$
est non vide. D'autre part, puisque $j = (t, s)$ est localement de type fini,
le morphisme $Z \to \Spec(k)$ est lui aussi localement de type fini (utiliser
Morphismes d'espaces algébriques, lemmes \ref{spaces-morphisms-lemma-permanence-finite-type}
et \ref{spaces-morphisms-lemma-base-change-finite-type}).
Ainsi, $Z$ possède un point à valeurs dans $k$ d'après
Morphismes d'espaces algébriques, lemme
\ref{spaces-morphisms-lemma-locally-finite-type-surjective-geometric-points},
et l'on conclut que $(\overline{u}, \overline{u}') \in j(R(k))$,
comme voulu. Cela achève la démonstration du lemme.
\end{proof}

\begin{lemma}
\label{lemma-set-theoretic-equivalence-geometric}
Dans la situation de la définition \ref{definition-set-theoretic-equivalence},
les conditions suivantes sont équivalentes :
\begin{enumerate}
\item Le morphisme $j$ est une relation d'équivalence ensembliste.
\item Le morphisme $j$ est universellement injectif et
$j(|R|) \subset |U \times_B U|$ contient l'image de
$|j'|$ pour chacun des morphismes $j'$ de l'équation (\ref{equation-list}).
\item Pour tout corps algébriquement clos $k$ au-dessus de $B$ et de cardinal suffisamment grand,
l'application $j : R(k) \to U(k) \times U(k)$ est injective et
son image est une relation d'équivalence.
\end{enumerate}
Si $j$ est décent, localement séparé ou quasi-séparé,
ces conditions équivalent aussi à
\begin{enumerate}
\item[(4)] Pour tout corps algébriquement clos $k$ au-dessus de $B$,
l'application $j : R(k) \to U(k) \times U(k)$ est injective et son image
est une relation d'équivalence.
\end{enumerate}
\end{lemma}

\begin{proof}
Les implications (1) $\Rightarrow$ (2) et (2) $\Rightarrow$ (3) résultent du
lemme \ref{lemma-set-theoretic-pre-equivalence-geometric}
et des définitions. Le même lemme montre que (3) implique que
$j$ est une prérelation d'équivalence ensembliste. Mais la condition
(3) implique aussi, bien entendu, que $j$ est universellement injectif ; voir
Morphismes d'espaces algébriques, lemme \ref{spaces-morphisms-lemma-universally-injective} ;
ainsi, $j$ est bien une relation d'équivalence ensembliste.
Nous savons donc que (1), (2) et (3) sont toutes équivalentes.

\medskip\noindent
La condition (4) implique (3) sans autre hypothèse sur $j$. Supposons que $j$
soit décent, localement séparé ou quasi-séparé, et que les conditions équivalentes
(1), (2), (3) soient satisfaites. D'après
Compléments sur les morphismes d'espaces algébriques,
lemme \ref{spaces-more-morphisms-lemma-when-universally-injective-radicial},
$j$ est radiciel.
Soit $k$ un corps algébriquement clos quelconque au-dessus de $B$. Soient
$\overline{u}, \overline{u}' \in U(k)$ tels que
$\overline{u} \sim_R \overline{u}'$. On voit que
$R \times_{U \times_B U, (\overline{u}, \overline{u}')} \Spec(k)$
est non vide. Comme $j$ est radiciel, sa réduction est donc le spectre d'un
corps purement inséparable sur $k$. Puisque $k = \overline{k}$, on voit que
c'est le spectre de $k$. Il existe donc un point $\overline{r} \in R(k)$
tel que $t(\overline{r}) = \overline{u}$ et $s(\overline{r}) = \overline{u}'$,
comme voulu.
\end{proof}

\begin{lemma}
\label{lemma-set-theoretic-equivalence}
Soit $S$ un schéma et soit $B$ un espace algébrique sur $S$.
Soit $j : R \to U \times_B U$ une prérelation sur $B$.
\begin{enumerate}
\item Si $j$ est une prérelation d'équivalence, alors $j$ est une
prérelation d'équivalence ensembliste. C'est en particulier le cas
lorsque $j$ provient d'un groupoïde en espaces algébriques ou de
l'action d'un espace algébrique en groupes sur $U$.
\item Si $j$ est une relation d'équivalence, alors $j$ est une
relation d'équivalence ensembliste.
\end{enumerate}
\end{lemma}

\begin{proof}
Omise.
\end{proof}

\begin{lemma}
\label{lemma-separates-orbits}
Soit $B \to S$ comme dans la section \ref{section-conventions-notation}.
Soit $j : R \to U \times_B U$ une prérelation.
Soit $\phi : U \to X$ un morphisme d'espaces algébriques sur $B$.
Considérons le diagramme
$$
\xymatrix{
(U \times_X U) \times_{(U \times_B U)} R \ar[d]^q \ar[r]_-p & R \ar[d]^j \\
U \times_X U \ar[r]^c & U \times_B U
}
$$
Alors :
\begin{enumerate}
\item Le morphisme $\phi$ est ensemblistement invariant si et seulement
si $p$ est surjectif.
\item Si $j$ est une prérelation d'équivalence ensembliste, alors
$\phi$ sépare les orbites si et seulement si $p$ et $q$ sont surjectifs.
\item Si $p$ et $q$ sont surjectifs, alors $j$ est une prérelation
d'équivalence ensembliste (et $\phi$ sépare les orbites).
\item Si $\phi$ est $R$-invariant et si $j$ est une prérelation d'équivalence
ensembliste, alors $\phi$ sépare les orbites si et seulement si le morphisme induit
$R \to U \times_X U$ est surjectif.
\end{enumerate}
\end{lemma}

\begin{proof}
Supposons que $\phi$ soit ensemblistement invariant. Cela signifie que, pour tout
corps algébriquement clos $k$ au-dessus de $B$ et tout $\overline{r} \in R(k)$,
on a $\phi(s(\overline{r})) = \phi(t(\overline{r}))$. Ainsi,
$((\phi(t(\overline{r})), \phi(s(\overline{r}))), \overline{r})$
définit un point du produit fibré qui s'envoie sur $\overline{r}$ par
$p$. Cela montre que $p$ est surjectif. Réciproquement, supposons $p$
surjectif. Choisissons $\overline{r} \in R(k)$. Puisque $p$ est surjectif, on
peut trouver une extension de corps $K/k$ et un point à valeurs dans $K$
$\tilde r$ du produit fibré tel que $p(\tilde r) = \overline{r}$.
Alors $q(\tilde r) \in U \times_X U$ s'envoie sur
$(t(\overline{r}), s(\overline{r}))$ dans $U \times_B U$, et l'on conclut
que $\phi(s(\overline{r})) = \phi(t(\overline{r}))$. Cela prouve
que $\phi$ est ensemblistement invariant.

\medskip\noindent
Les démonstrations de (2), (3) et (4) sont omises. Indication : supposer que $k$ est un
corps algébriquement clos au-dessus de $B$, de grand cardinal. Considérer le
diagramme d'ensembles associé
$$
\xymatrix{
(U(k) \times_{X(k)} U(k)) \times_{U(k) \times U(k)} R(k) \ar[d]^q \ar[r]_-p &
R(k) \ar[d]^j \\
U(k) \times_{X(k)} U(k) \ar[r]^c & U(k) \times U(k)
}
$$
Grâce aux lemmes précédents, les équivalences formulées en (2), (3) et (4) deviennent
des questions ensemblistes relatives au diagramme que l'on vient d'afficher, en utilisant
le fait que la surjectivité se traduit par la surjectivité sur les points à valeurs dans $k$, d'après
Morphismes d'espaces algébriques, lemme \ref{spaces-morphisms-lemma-large-enough}.
\end{proof}

\noindent
Comme nous avons vu plus haut que la notion de morphisme ensemblistement
invariant est assez faible dans la catégorie des espaces algébriques,
nous définissons comme suit un espace d'orbites pour une prérelation.

\begin{definition}
\label{definition-orbit-space}
Soit $B \to S$ comme dans la section \ref{section-conventions-notation}.
Soit $j : R \to U \times_B U$ une prérelation.
On dit que $\phi : U \to X$ est un {\it espace d'orbites pour $R$} si
\begin{enumerate}
\item $\phi$ est $R$-invariant ;
\item $\phi$ sépare les $R$-orbites ; et
\item $\phi$ est surjectif.
\end{enumerate}
\end{definition}

\noindent
La définition de la séparation des $R$-orbites fait intervenir des
points à valeurs dans des corps algébriquement clos. Mais, comme on l'a vu,
cela correspond dans de nombreux cas à la surjectivité de certains
morphismes d'espaces algébriques canoniquement associés.
Résumons la discussion précédente dans la caractérisation suivante
des espaces d'orbites.

\begin{lemma}
\label{lemma-orbit-space}
Soit $B \to S$ comme dans la section \ref{section-conventions-notation}.
Soit $j : R \to U \times_B U$ une prérelation d'équivalence
ensembliste. Un morphisme $\phi : U \to X$ est un espace d'orbites pour $R$ si et seulement si
\begin{enumerate}
\item $\phi \circ s = \phi \circ t$, c'est-à-dire que $\phi$ est invariant ;
\item le morphisme induit $(t, s) : R \to U \times_X U$ est surjectif ; et
\item le morphisme $\phi : U \to X$ est surjectif.
\end{enumerate}
Cette caractérisation s'applique par exemple si $j$ est une prérelation d'équivalence,
s'il provient d'un groupoïde en espaces algébriques sur $B$, ou s'il provient de l'action
d'un espace algébrique en groupes sur $B$ sur $U$.
\end{lemma}

\begin{proof}
Cela résulte immédiatement du lemme \ref{lemma-separates-orbits}, partie (4).
\end{proof}

\noindent
Dans le lemme suivant, il ne suffit (probablement) pas de supposer seulement que
les morphismes $s, t$ sont localement de type fini. En effet,
il peut arriver qu'une application $\phi : U \to X$ soit un espace d'orbites sans être
localement de type fini. Dans ce cas, $U(k) \to X(k)$ peut ne pas
être surjective pour tout corps algébriquement clos $k$ au-dessus de $B$.

\begin{lemma}
\label{lemma-orbit-space-locally-finite-type-over-base}
Soit $B \to S$ comme dans la section \ref{section-conventions-notation}.
Soit $j = (t, s) : R \to U \times_B U$ une prérelation.
Supposons que $R, U$ soient localement de type fini sur $B$.
Soit $\phi : U \to X$ un morphisme $R$-invariant d'espaces algébriques sur $B$.
Alors $\phi$ est un espace d'orbites pour $R$ si et seulement si l'application naturelle
$$
U(k)/\big(\text{equivalence relation generated by }j(R(k))\big)
\longrightarrow
X(k)
$$
est bijective pour tout corps algébriquement clos $k$ au-dessus de $B$.
\end{lemma}

\begin{proof}
Remarquons que, puisque $U$, $R$ sont localement de type fini sur $B$, tous les
morphismes $s, t, j, \phi$ sont localement de type fini ; voir
Morphismes d'espaces algébriques, lemme \ref{spaces-morphisms-lemma-permanence-finite-type}.
Nous utiliserons aussi sans autre mention
Morphismes d'espaces algébriques, lemme
\ref{spaces-morphisms-lemma-locally-finite-type-surjective-geometric-points}.
Supposons que $\phi$ soit un espace d'orbites. Soit $k$ un corps algébriquement clos
quelconque au-dessus de $B$. Soit $\overline{x} \in X(k)$. Considérons
$U \times_{\phi, X, \overline{x}} \Spec(k)$.
C'est un espace algébrique non vide,
localement de type fini sur $k$. Il possède donc un point à valeurs dans $k$.
Cela montre que l'application affichée dans le lemme est surjective.
Supposons que $\overline{u}, \overline{u}' \in U(k)$ aient la même
image dans $X(k)$. D'après la
définition \ref{definition-set-theoretically-invariant},
cela signifie que $\overline{u}, \overline{u}'$ sont dans la même
$R$-orbite. D'après le lemme \ref{lemma-geometric-orbits}, cela signifie qu'ils
sont équivalents pour la relation d'équivalence engendrée par
$j(R(k))$. Ainsi, le morphisme affiché est injectif.

\medskip\noindent
Réciproquement, supposons que l'application affichée soit bijective pour tout corps
algébriquement clos $k$ au-dessus de $B$. Cette condition implique clairement que $\phi$
est surjectif. Nous avons déjà supposé que $\phi$ est $R$-invariant.
Enfin, l'injectivité de toutes les applications affichées implique que
$\phi$ sépare les orbites. Par conséquent, $\phi$ est un espace d'orbites.
\end{proof}










\section{Quotients grossiers}
\label{section-coarse}

\noindent
Nous ajoutons seulement cette notion ici afin de pouvoir dire plus tard que les quotients
grossiers correspondent aux espaces de modules grossiers (ou aux schémas de modules).

\begin{definition}
\label{definition-coarse}
Soit $S$ un schéma et soit $B$ un espace algébrique sur $S$.
Soit $j : R \to U \times_B U$ une prérelation.
Un morphisme $\phi : U \to X$ d'espaces algébriques sur $B$
est appelé {\it quotient grossier} si
\begin{enumerate}
\item $\phi$ est un quotient catégorique ; et
\item $\phi$ est un espace d'orbites.
\end{enumerate}
Si $S = B$ et si $U$, $R$ sont tous deux des schémas, on dit qu'un morphisme de schémas
$\phi : U \to X$ est un {\it quotient grossier en schémas} si
\begin{enumerate}
\item $\phi$ est un quotient catégorique en schémas ; et
\item $\phi$ est un espace d'orbites.
\end{enumerate}
\end{definition}

\noindent
Dans de nombreuses situations, les espaces algébriques $R$ et $U$ sont localement de type fini
sur $B$, et la condition d'espace d'orbites signifie simplement que
$$
U(k)/\big(\text{equivalence relation generated by }j(R(k))\big)
\cong
X(k)
$$
pour tout corps algébriquement clos $k$. Voir le
lemme \ref{lemma-orbit-space-locally-finite-type-over-base}.
Si $j$ est en outre une prérelation d'équivalence (ensembliste), alors la condition
équivaut simplement à ce que $U(k)/j(R(k)) \to X(k)$ soit bijective pour tout
corps algébriquement clos $k$.









\section{Propriétés topologiques}
\label{section-topological}

\noindent
Soit $S$ un schéma et soit $B$ un espace algébrique sur $S$.
Soit $j : R \to U \times_B U$ une prérelation.
On dit qu'un sous-ensemble $T \subset |U|$ est {\it $R$-invariant} si
$s^{-1}(T) = t^{-1}(T)$ comme sous-ensembles de $|R|$.
Remarquons que, si $T$ est fermé, il peut arriver que
le sous-espace fermé réduit correspondant de $U$ ne soit pas $R$-invariant
(au sens de
Groupoïdes in Espaces, définition
\ref{spaces-groupoids-definition-invariant-open}),
car les images réciproques $s^{-1}(T)$, $t^{-1}(T)$ peuvent ne pas être réduites.
Voici quelques conditions que l'on peut considérer pour un
morphisme invariant $\phi : U \to X$.

\begin{definition}
\label{definition-topological}
Soit $S$ un schéma et soit $B$ un espace algébrique sur $S$.
Soit $j : R \to U \times_B U$ une prérelation.
Soit $\phi : U \to X$ un morphisme $R$-invariant d'espaces algébriques sur $B$.
\begin{enumerate}
\item
\label{item-submersive}
Le morphisme $\phi$ est submersif.
\item
\label{item-invariant-closed}
Pour tout sous-ensemble fermé $R$-invariant $Z \subset |U|$, l'image
$\phi(Z)$ est fermée dans $|X|$.
\item
\label{item-intersect-invariant-closed}
La condition (\ref{item-invariant-closed}) est satisfaite et, pour toute paire de
sous-ensembles fermés $R$-invariants $Z_1, Z_2 \subset |U|$, on a
$$
\phi(Z_1 \cap Z_2) = \phi(Z_1) \cap \phi(Z_2)
$$
\item Le morphisme $(t, s) : R \to U \times_X U$ est universellement submersif.
\label{item-strong}
\end{enumerate}
Pour chacune de ces propriétés, on peut aussi exiger qu'elle subsiste après tout
changement de base plat, ou après tout changement de base ; voir la
définition \ref{definition-base-change}. On dit alors que la condition
(\ref{item-submersive}),
(\ref{item-invariant-closed}),
(\ref{item-intersect-invariant-closed}) ou
(\ref{item-strong}) est satisfaite {\it uniformément} ou {\it universellement}.
\end{definition}









\section{Fonctions invariantes}
\label{section-functions}

\noindent
Dans certains cas, il est commode de déterminer précisément le faisceau structural
d'un quotient en exigeant que toute fonction invariante soit une section locale
du faisceau structural du quotient.

\begin{definition}
\label{definition-functions}
Soit $S$ un schéma et soit $B$ un espace algébrique sur $S$.
Soit $j : R \to U \times_B U$ une prérelation.
Soit $\phi : U \to X$ un morphisme $R$-invariant.
Notons $\phi' = \phi \circ s = \phi \circ t : R \to X$.
\begin{enumerate}
\item On note $(\phi_*\mathcal{O}_U)^R$ la sous-$\mathcal{O}_X$-algèbre
de $\phi_*\mathcal{O}_U$ qui est le noyau de la double flèche
$$
\xymatrix{
\phi_*\mathcal{O}_U
\ar@<1ex>[rr]^{\phi_*s^\sharp}
\ar@<-1ex>[rr]_{\phi_*t^\sharp}
& &
\phi'_*\mathcal{O}_R
}
$$
sur $X_\etale$. On l'appelle parfois le
{\it faisceau des fonctions $R$-invariantes sur $X$}.
\item On dit que {\it les fonctions sur $X$ sont les fonctions $R$-invariantes sur
$U$} si l'application naturelle $\mathcal{O}_X \to (\phi_*\mathcal{O}_U)^R$
est un isomorphisme.
\end{enumerate}
\end{definition}

\noindent
On peut bien sûr exiger que cette propriété subsiste après tout
changement de base (plat ou quelconque), ce qui donne une notion (uniforme ou) universelle.
Cette condition est souvent ajoutée
à d'autres afin d'obtenir un quotient davantage déterminé. Une grande partie de
la motivation de toute la théorie provient bien sûr du cas particulier suivant :
$U = \Spec(A)$ est un schéma affine sur un corps
$S = B = \Spec(k)$ et $R = G \times U$,
où $G$ est un schéma en groupes affine sur $k$. Dans ce cas,
on peut choisir comme quotient
$$
X = \Spec(A^G)
$$
de sorte qu'au moins la condition de la définition précédente soit satisfaite.
Bien que cette construction soit agréable, elle ne donne souvent pas le bon
quotient ; par exemple, si $U = \text{GL}_{n, k}$ et si $G$ est le groupe des
matrices triangulaires supérieures, elle donne $X = \Spec(k)$, alors qu'il
existe un bien meilleur quotient (à savoir la variété de drapeaux).









\section{Bons quotients}
\label{section-good}

\noindent
La définition suivante est particulièrement utile lorsqu'on prend des quotients
par des actions de groupes.

\begin{definition}
\label{definition-good}
Soit $S$ un schéma et soit $B$ un espace algébrique sur $S$.
Soit $j : R \to U \times_B U$ une prérelation.
Un morphisme $\phi : U \to X$ d'espaces algébriques sur $B$
est appelé {\it bon quotient} si
\begin{enumerate}
\item $\phi$ est invariant ;
\item $\phi$ est affine ;
\item $\phi$ est surjectif ;
\item la condition (\ref{item-intersect-invariant-closed}) est satisfaite universellement ; et
\item les fonctions sur $X$ sont les fonctions $R$-invariantes sur $U$.
\end{enumerate}
\end{definition}

\noindent
Dans \cite{seshadri_quotients}, Seshadri donne presque la même définition,
mais, à la place de (4), il exige simplement que la condition
(\ref{item-intersect-invariant-closed}) soit satisfaite, sans exiger
qu'elle le soit universellement.








\section{Quotients géométriques}
\label{section-geometric}

\noindent
C'est la définition de Mumford d'un quotient géométrique (du moins celle
de la première édition de GIT ; pour autant que nous puissions en juger, les éditions ultérieures
ont remplacé « universellement submersif » par « submersif »).

\begin{definition}
\label{definition-geometric}
Soit $S$ un schéma et soit $B$ un espace algébrique sur $S$.
Soit $j : R \to U \times_B U$ une prérelation.
Un morphisme $\phi : U \to X$ d'espaces algébriques sur $B$
est appelé {\it quotient géométrique} si
\begin{enumerate}
\item $\phi$ est un espace d'orbites ;
\item la condition (\ref{item-submersive}) est satisfaite universellement, c'est-à-dire
que $\phi$ est universellement submersif ; et
\item les fonctions sur $X$ sont les fonctions $R$-invariantes sur $U$.
\end{enumerate}
\end{definition}










\begin{multicols}{2}[\section{Autres chapitres}]
\noindent
Préliminaires
\begin{enumerate}
\item \hyperref[introduction-section-phantom]{Introduction}
\item \hyperref[conventions-section-phantom]{Conventions}
\item \hyperref[sets-section-phantom]{Théorie des ensembles}
\item \hyperref[categories-section-phantom]{Catégories}
\item \hyperref[topology-section-phantom]{Topologie}
\item \hyperref[sheaves-section-phantom]{Faisceaux sur les espaces}
\item \hyperref[sites-section-phantom]{Sites et faisceaux}
\item \hyperref[stacks-section-phantom]{Champs}
\item \hyperref[fields-section-phantom]{Corps}
\item \hyperref[algebra-section-phantom]{Algèbre commutative}
\item \hyperref[brauer-section-phantom]{Groupes de Brauer}
\item \hyperref[homology-section-phantom]{Algèbre homologique}
\item \hyperref[derived-section-phantom]{Catégories dérivées}
\item \hyperref[simplicial-section-phantom]{Méthodes simpliciales}
\item \hyperref[more-algebra-section-phantom]{Compléments d'algèbre}
\item \hyperref[smoothing-section-phantom]{Lissage des morphismes d'anneaux}
\item \hyperref[modules-section-phantom]{Faisceaux de modules}
\item \hyperref[sites-modules-section-phantom]{Modules sur les sites}
\item \hyperref[injectives-section-phantom]{Objets injectifs}
\item \hyperref[cohomology-section-phantom]{Cohomologie des faisceaux}
\item \hyperref[sites-cohomology-section-phantom]{Cohomologie sur les sites}
\item \hyperref[dga-section-phantom]{Algèbre différentielle graduée}
\item \hyperref[dpa-section-phantom]{Algèbre à puissances divisées}
\item \hyperref[sdga-section-phantom]{Faisceaux différentiels gradués}
\item \hyperref[hypercovering-section-phantom]{Hyperrecouvrements}
\end{enumerate}
Schémas
\begin{enumerate}
\setcounter{enumi}{25}
\item \hyperref[schemes-section-phantom]{Schémas}
\item \hyperref[constructions-section-phantom]{Constructions de schémas}
\item \hyperref[properties-section-phantom]{Propriétés des schémas}
\item \hyperref[morphisms-section-phantom]{Morphismes de schémas}
\item \hyperref[coherent-section-phantom]{Cohomologie des schémas}
\item \hyperref[divisors-section-phantom]{Diviseurs}
\item \hyperref[limits-section-phantom]{Limites de schémas}
\item \hyperref[varieties-section-phantom]{Variétés}
\item \hyperref[topologies-section-phantom]{Topologies sur les schémas}
\item \hyperref[descent-section-phantom]{Descente}
\item \hyperref[perfect-section-phantom]{Catégories dérivées de schémas}
\item \hyperref[more-morphisms-section-phantom]{Compléments sur les morphismes}
\item \hyperref[flat-section-phantom]{Compléments sur la platitude}
\item \hyperref[groupoids-section-phantom]{Schémas en groupoïdes}
\item \hyperref[more-groupoids-section-phantom]{Compléments sur les schémas en groupoïdes}
\item \hyperref[etale-section-phantom]{Morphismes étales de schémas}
\end{enumerate}
Thèmes de théorie des schémas
\begin{enumerate}
\setcounter{enumi}{41}
\item \hyperref[chow-section-phantom]{Homologie de Chow}
\item \hyperref[intersection-section-phantom]{Théorie de l'intersection}
\item \hyperref[pic-section-phantom]{Schémas de Picard des courbes}
\item \hyperref[weil-section-phantom]{Théories cohomologiques de Weil}
\item \hyperref[adequate-section-phantom]{Modules adéquats}
\item \hyperref[dualizing-section-phantom]{Complexes dualisants}
\item \hyperref[duality-section-phantom]{Dualité pour les schémas}
\item \hyperref[discriminant-section-phantom]{Discriminants et différentes}
\item \hyperref[derham-section-phantom]{Cohomologie de de Rham}
\item \hyperref[local-cohomology-section-phantom]{Cohomologie locale}
\item \hyperref[algebraization-section-phantom]{Géométrie algébrique et formelle}
\item \hyperref[curves-section-phantom]{Courbes algébriques}
\item \hyperref[resolve-section-phantom]{Résolution des surfaces}
\item \hyperref[models-section-phantom]{Réduction semi-stable}
\item \hyperref[functors-section-phantom]{Foncteurs et morphismes}
\item \hyperref[equiv-section-phantom]{Catégories dérivées de variétés}
\item \hyperref[pione-section-phantom]{Groupes fondamentaux de schémas}
\item \hyperref[etale-cohomology-section-phantom]{Cohomologie étale}
\item \hyperref[crystalline-section-phantom]{Cohomologie cristalline}
\item \hyperref[proetale-section-phantom]{Cohomologie pro-étale}
\item \hyperref[relative-cycles-section-phantom]{Cycles relatifs}
\item \hyperref[more-etale-section-phantom]{Compléments de cohomologie étale}
\item \hyperref[trace-section-phantom]{La formule des traces}
\end{enumerate}
Espaces algébriques
\begin{enumerate}
\setcounter{enumi}{64}
\item \hyperref[spaces-section-phantom]{Espaces algébriques}
\item \hyperref[spaces-properties-section-phantom]{Propriétés des espaces algébriques}
\item \hyperref[spaces-morphisms-section-phantom]{Morphismes d'espaces algébriques}
\item \hyperref[decent-spaces-section-phantom]{Espaces algébriques décents}
\item \hyperref[spaces-cohomology-section-phantom]{Cohomologie des espaces algébriques}
\item \hyperref[spaces-limits-section-phantom]{Limites d'espaces algébriques}
\item \hyperref[spaces-divisors-section-phantom]{Diviseurs sur les espaces algébriques}
\item \hyperref[spaces-over-fields-section-phantom]{Espaces algébriques sur des corps}
\item \hyperref[spaces-topologies-section-phantom]{Topologies sur les espaces algébriques}
\item \hyperref[spaces-descent-section-phantom]{Descente et espaces algébriques}
\item \hyperref[spaces-perfect-section-phantom]{Catégories dérivées d'espaces}
\item \hyperref[spaces-more-morphisms-section-phantom]{Compléments sur les morphismes d'espaces}
\item \hyperref[spaces-flat-section-phantom]{Platitude sur les espaces algébriques}
\item \hyperref[spaces-groupoids-section-phantom]{Groupoïdes dans les espaces algébriques}
\item \hyperref[spaces-more-groupoids-section-phantom]{Compléments sur les groupoïdes dans les espaces}
\item \hyperref[bootstrap-section-phantom]{Amorçage}
\item \hyperref[spaces-pushouts-section-phantom]{Sommes amalgamées d'espaces algébriques}
\end{enumerate}
Thèmes de géométrie
\begin{enumerate}
\setcounter{enumi}{81}
\item \hyperref[spaces-chow-section-phantom]{Groupes de Chow des espaces}
\item \hyperref[groupoids-quotients-section-phantom]{Quotients de groupoïdes}
\item \hyperref[spaces-more-cohomology-section-phantom]{Compléments sur la cohomologie des espaces}
\item \hyperref[spaces-simplicial-section-phantom]{Espaces simpliciaux}
\item \hyperref[spaces-duality-section-phantom]{Dualité pour les espaces}
\item \hyperref[formal-spaces-section-phantom]{Espaces algébriques formels}
\item \hyperref[restricted-section-phantom]{Algébrisation des espaces formels}
\item \hyperref[spaces-resolve-section-phantom]{Résolution des surfaces revisitée}
\end{enumerate}
Théorie des déformations
\begin{enumerate}
\setcounter{enumi}{89}
\item \hyperref[formal-defos-section-phantom]{Théorie formelle des déformations}
\item \hyperref[defos-section-phantom]{Théorie des déformations}
\item \hyperref[cotangent-section-phantom]{Le complexe cotangent}
\item \hyperref[examples-defos-section-phantom]{Problèmes de déformation}
\end{enumerate}
Champs algébriques
\begin{enumerate}
\setcounter{enumi}{93}
\item \hyperref[algebraic-section-phantom]{Champs algébriques}
\item \hyperref[examples-stacks-section-phantom]{Exemples de champs}
\item \hyperref[stacks-sheaves-section-phantom]{Faisceaux sur les champs algébriques}
\item \hyperref[criteria-section-phantom]{Critères de représentabilité}
\item \hyperref[artin-section-phantom]{Axiomes d'Artin}
\item \hyperref[quot-section-phantom]{Espaces de Quot et de Hilbert}
\item \hyperref[stacks-properties-section-phantom]{Propriétés des champs algébriques}
\item \hyperref[stacks-morphisms-section-phantom]{Morphismes de champs algébriques}
\item \hyperref[stacks-limits-section-phantom]{Limites de champs algébriques}
\item \hyperref[stacks-cohomology-section-phantom]{Cohomologie des champs algébriques}
\item \hyperref[stacks-perfect-section-phantom]{Catégories dérivées de champs}
\item \hyperref[stacks-introduction-section-phantom]{Introduction aux champs algébriques}
\item \hyperref[stacks-more-morphisms-section-phantom]{Compléments sur les morphismes de champs}
\item \hyperref[stacks-geometry-section-phantom]{La géométrie des champs}
\end{enumerate}
Thèmes de théorie des espaces de modules
\begin{enumerate}
\setcounter{enumi}{107}
\item \hyperref[moduli-section-phantom]{Champs de modules}
\item \hyperref[moduli-curves-section-phantom]{Espaces de modules de courbes}
\end{enumerate}
Divers
\begin{enumerate}
\setcounter{enumi}{109}
\item \hyperref[examples-section-phantom]{Exemples}
\item \hyperref[exercises-section-phantom]{Exercices}
\item \hyperref[guide-section-phantom]{Guide bibliographique}
\item \hyperref[desirables-section-phantom]{Desiderata}
\item \hyperref[coding-section-phantom]{Style de programmation}
\item \hyperref[obsolete-section-phantom]{Obsolète}
\item \hyperref[fdl-section-phantom]{Licence de documentation libre GNU}
\item \hyperref[index-section-phantom]{Index généré automatiquement}
\end{enumerate}
\end{multicols}


\bibliography{my}
\bibliographystyle{amsalpha}

\end{document}

